\section{Conclusion}
\label{sec:conclusion}

We introduced the DRM Transformer, a decoder-only Transformer that replaces
dot-product attention with geodesic attention over a learned Directional
Relational Manifold. The architecture combines a low-rank MetricNet,
gravitational deformation, relativistic gamma-scaling, and variable
dimensionality. Based on the DRM and relativistic dynamics papers, we
implemented a trainable neural version of emergent dimensionality and toroidal
structure.

The most important result so far is the observation of a seed-robust
homological signature compatible with a near-toroidal closed surface in the
small\_3.5M baseline. The controls suggest that training moves the DRM manifold
from an initial high-genus unstable topology toward this signature: in two out
of three trained seeds, the median topology matches $H_1=2$, $H_2=1$, the
expected signature of $T^2$; all trained seeds show a \emph{Near T2} signal,
while random initialization does not. This does not complete validation of the
method, since the strict criterion $f_{T^2}\geq 0.60$ has not been reached, but
it demonstrates that DRM geometry can be operationalized in language models and
evaluated with objective topological tools.

Future work includes scaling the model, testing ablations at 50M, 350M, and
1B+ parameters, comparing against standard attention under equivalent compute
budgets, and studying the relation between DRM topology and perplexity. A
central validation point will be comparing four regimes: explicit torus, torus
with annealing to zero, generic geometry without toroidal target, and absence
of toroidal regularization. If regimes without a toroidal target maintain
$H_1=2$, $H_2=1$ with higher stability at scale, evidence for topological
emergence will be substantially stronger than in the explicitly prescriptive
regime.
